\documentclass{article}

\usepackage{mathtools}
\usepackage{amssymb}
\usepackage{nicematrix}
\usepackage{zref-clever}
\usepackage{../latex-course}

% Paired delimiters declared here
\DeclarePairedDelimiter{\abs}{\lvert}{\rvert}
\DeclarePairedDelimiterX{\set}[2]{\lbrace}{\rbrace}{#1 \;\delimsize\vert\; \mathopen{} #2}
\DeclarePairedDelimiterX{\augment}[2]{\lbrack}{\rbrack}{#1 \;\delimsize\vert\; \mathopen{} #2}

% Operator names declared here
\DeclareMathOperator{\rank}{rank}

\courseweek{04}{Math Formatting -- Part 2}
\coursestudents{Sandrine Pudukkattukaran Joy \and Vincent Kamusella}

\title{Linear Systems of Equations}
\author{Sandrine Pudukkattukaran Joy (4743017), Vincent Kamusella (4752668)}
\begin{document}
\maketitle

\section{A Motivating Example}
\label{sec:example}
Over the field $\mathbb{R}$, consider the linear system
\begin{subequations}
	\label{eq:system}
	\begin{align}
		2x + y & = 5, \label{eq:equation:sum}\\
		x - y & = 1. \label{eq:equation:diff}
	\end{align}
\end{subequations}
Adding the two equations eliminates $y$:
\begin{equation}
	\label{eq:x}
	\tag{\ref{eq:equation:sum}+\ref{eq:equation:diff}}
	3x = 6.
\end{equation}
From \eqref{eq:x}, we get $x = 2$, and substituting into \eqref{eq:equation:diff} (or \eqref{eq:equation:sum}) then gives $y = 1$.
The unique solution is therefore
\begin{equation*}
	\begin{bmatrix}
		x \\ y
	\end{bmatrix}
	=
	\begin{bmatrix}
		2 \\ 1
	\end{bmatrix}.
\end{equation*}

\section{Linear Systems in General}
\label{sec:general}
Let $K$ be a field and $m, n \in \mathbb{N}$.
A \emph{linear system} of $n$ equations in $m$ unknowns is an equation $A x = b$ with coefficient matrix $A \in K^{n \times m}$, right-hand side $b \in K^n$ and variable $x \in K^m$:
\[
	A x
	=
	\begin{bNiceMatrix}
		a_{11} & \Cdots & a_{1m} \\
		\Vdots & \Ddots & \Vdots \\
		a_{n1} & \Cdots & a_{nm}
	\end{bNiceMatrix}
	\begin{bNiceMatrix}
		x_1 \\ \Vdots \\ x_m
	\end{bNiceMatrix}
	=
	\begin{bNiceMatrix}
		b_1 \\ \Vdots \\ b_n
	\end{bNiceMatrix}
	=
	b.
\]
We seek its \emph{solution set}
\begin{equation}
	\label{eq:solution-set}
	L(A, b) \coloneqq \set{x \in K^m}{A x = b}.
\end{equation}
The system is solvable if and only if $b$ lies in the column space of $A$, which is equivalent to $\rank(A) = \rank(\augment{A}{b})$, where $\augment{A}{b}\in K^{n \times (m + 1)}$ is the right hand side augmented coefficient matrix.
Provided we find any one particular solution $x_0$, then $L(A, b) = x_0 + L(A, 0)$, and $L(A, 0)$ is a subspace of $K^m$ of dimension $m - \rank(A)$.
Accordingly, the set defined in \eqref{eq:solution-set} is either empty or has the same cardinality as the solution set to the homogeneous system.
More specifically, we have that
\begin{equation*}
\abs{L(A, b)} =
\begin{cases}
	0 & \text{if } \rank(A) \neq \rank(\augment{A}{b}), \\
	\abs{L(A, 0)} = 1 & \text{if } \rank(A) = \rank(\augment{A}{b}) = m, \\
	\abs{L(A, 0)} > 1 & \text{if } \rank(A) = \rank(\augment{A}{b}) < m.
\end{cases}
\end{equation*}
For the system in \eqref{eq:system}, we have $n = m = 2$ and $\rank\begin{psmallmatrix} 2 & 1 \\ 1 & -1 \end{psmallmatrix} =\rank\begin{pNiceArray}{cc|c}[small] 2 & 1 & 5 \\ 1 & -1 & 1 \end{pNiceArray} = m$, so it is uniquely solvable, in agreement with the explanation in \zcref{sec:example}.
\end{document}
